% --- Archon physics formalization source begin ---
% archon:physics
% archon:covers PhyXMiniProblems/problem_phyx_mini_0215.lean
% archon:source-report reports/phyx_mini/problem_phyx_mini_0215.source.json

\chapter{Physics problem phyx\_mini\_0215}
\label{ch:PhyXMiniProblems_problem_phyx_mini_0215}

\paragraph{Problem source.}
\#\# Physical scenario

An ocean fishing boat is drifting just above a school of tuna on a foggy day. Without warning, an engine backfire occurs on another boat \textbackslash{}( 1.55\textbackslash{},\textbackslash{}mathrm\{km\} \textbackslash{}) away (figure).

\#\# Question

How much time elapses before the backfire is heard by the fishermen?

\#\# Answer choices

- A: A: 3.52s

- B: B: 6.52s

- C: C: 5.52s

- D: D: 4.52s

\#\# Figure caption (auxiliary; use the image as primary evidence)

The image depicts a seascape with two fishing boats and a group of fish. 

1. **Foreground:**
   - A group of four fish, labeled as (a), is shown swimming in the water.

2. **Middle Ground:**
   - A fishing boat with an orange bottom and white top is positioned to the left, near a lighthouse. This boat is labeled as (b).

3. **Background:**
   - Another fishing boat, with a green bottom and white top, is situated further away on the right side of the image.

4. **Annotations:**
   - A distance measurement of "1.55 km" is indicated between the two boats, showing the distance from the larger boat (b) to the smaller boat in the background.

5. **Other Details:**
   - The sky is depicted as partly cloudy.

\#\# Dataset metadata

Sound; Physical Model Grounding Reasoning, Predictive Reasoning

\paragraph{Recorded answer/context.}
D: D: 4.52s

\paragraph{Figure/image path.}
/root/proposal\_for\_physic/science-mango/phyx\_mini\_run/phyx\_data/test\_image/215.png

\paragraph{Formalization target.}
create a compiling Lean file with sorry bodies at `PhyXMiniProblems/problem\_phyx\_mini\_0215.lean`. The Lean declarations must preserve the physical quantities, dimensions or dimensional roles, figure labels, governing-law hypotheses, and final relation expressed by this problem.
Use Mathlib/Physlib names found through LeanExplore where available. If a physics API is missing, introduce faithful local abstractions rather than scalar placeholder aliases.

\begin{theorem}[Physics formalization target]
\label{thm:physics:phyx_mini_0215:target}
\lean{PhyXMiniProblems.ProblemPhyXMini0215.backfire_is_heard_after_answer_D}
\uses{def:physics:phyx-mini-0215:phyxminiproblems-problemphyxmini0215-durationinseconds, def:physics:phyx-mini-0215:phyxminiproblems-problemphyxmini0215-boatbackfiresetup, def:physics:phyx-mini-0215:phyxminiproblems-problemphyxmini0215-matchesproblemandfigure, def:physics:phyx-mini-0215:phyxminiproblems-problemphyxmini0215-hasphysicalparameters, def:physics:phyx-mini-0215:phyxminiproblems-problemphyxmini0215-usesstandardairsoundspeed, def:physics:phyx-mini-0215:phyxminiproblems-problemphyxmini0215-obeysconstantspeedsoundpropagation, def:physics:phyx-mini-0215:phyxminiproblems-problemphyxmini0215-matchesanswerchoice}
The assigned autoformalize agent should translate this physics problem into Lean declarations in the covered file, with theorem and lemma proof bodies written as `by sorry`.
\end{theorem}
\begin{proof}
This is an autoformalization task, not a proof task. Produce statements that can later be proved by the physics prover without weakening the physical model.
\end{proof}

% --- Archon declaration topology begin ---
\section{Declaration topology}
\begin{definition}[Length Quantity]
  \label{def:physics:phyx-mini-0215:phyxminiproblems-problemphyxmini0215-lengthquantity}
  \lean{PhyXMiniProblems.ProblemPhyXMini0215.LengthQuantity}
  A real-valued physical length, independent of the chosen unit readout
\end{definition}
\begin{proof}
  This is the declared model definition of Length Quantity.
\end{proof}
\begin{definition}[Duration Quantity]
  \label{def:physics:phyx-mini-0215:phyxminiproblems-problemphyxmini0215-durationquantity}
  \lean{PhyXMiniProblems.ProblemPhyXMini0215.DurationQuantity}
  A real-valued physical duration, independent of the chosen unit readout
\end{definition}
\begin{proof}
  This is the declared model definition of Duration Quantity.
\end{proof}
\begin{definition}[Speed Quantity]
  \label{def:physics:phyx-mini-0215:phyxminiproblems-problemphyxmini0215-speedquantity}
  \lean{PhyXMiniProblems.ProblemPhyXMini0215.SpeedQuantity}
  A real-valued physical speed, carrying length-per-time dimension
\end{definition}
\begin{proof}
  This is the declared model definition of Speed Quantity.
\end{proof}
\begin{definition}[length Readout]
  \label{def:physics:phyx-mini-0215:phyxminiproblems-problemphyxmini0215-lengthreadout}
  \lean{PhyXMiniProblems.ProblemPhyXMini0215.lengthReadout}
  \uses{def:physics:phyx-mini-0215:phyxminiproblems-problemphyxmini0215-lengthquantity}
  Read a physical length as a real number in the selected length unit
\end{definition}
\begin{proof}
  This is the declared model definition of length Readout.
\end{proof}
\begin{definition}[duration Readout]
  \label{def:physics:phyx-mini-0215:phyxminiproblems-problemphyxmini0215-durationreadout}
  \lean{PhyXMiniProblems.ProblemPhyXMini0215.durationReadout}
  \uses{def:physics:phyx-mini-0215:phyxminiproblems-problemphyxmini0215-durationquantity}
  Read a physical duration as a real number in the selected time unit
\end{definition}
\begin{proof}
  This is the declared model definition of duration Readout.
\end{proof}
\begin{definition}[speed In Meters Per Second]
  \label{def:physics:phyx-mini-0215:phyxminiproblems-problemphyxmini0215-speedinmeterspersecond}
  \lean{PhyXMiniProblems.ProblemPhyXMini0215.speedInMetersPerSecond}
  \uses{def:physics:phyx-mini-0215:phyxminiproblems-problemphyxmini0215-speedquantity}
  Read a physical speed in metres per second
\end{definition}
\begin{proof}
  This is the declared model definition of speed In Meters Per Second.
\end{proof}
\begin{definition}[length In Kilometers]
  \label{def:physics:phyx-mini-0215:phyxminiproblems-problemphyxmini0215-lengthinkilometers}
  \lean{PhyXMiniProblems.ProblemPhyXMini0215.lengthInKilometers}
  \uses{def:physics:phyx-mini-0215:phyxminiproblems-problemphyxmini0215-lengthquantity, def:physics:phyx-mini-0215:phyxminiproblems-problemphyxmini0215-lengthreadout}
  Kilometre readout used by the 1.55 km figure annotation
\end{definition}
\begin{proof}
  This is the declared model definition of length In Kilometers.
\end{proof}
\begin{definition}[length In Meters]
  \label{def:physics:phyx-mini-0215:phyxminiproblems-problemphyxmini0215-lengthinmeters}
  \lean{PhyXMiniProblems.ProblemPhyXMini0215.lengthInMeters}
  \uses{def:physics:phyx-mini-0215:phyxminiproblems-problemphyxmini0215-lengthquantity, def:physics:phyx-mini-0215:phyxminiproblems-problemphyxmini0215-lengthreadout}
  Metre readout used with the standard sound speed
\end{definition}
\begin{proof}
  This is the declared model definition of length In Meters.
\end{proof}
\begin{definition}[duration In Seconds]
  \label{def:physics:phyx-mini-0215:phyxminiproblems-problemphyxmini0215-durationinseconds}
  \lean{PhyXMiniProblems.ProblemPhyXMini0215.durationInSeconds}
  \uses{def:physics:phyx-mini-0215:phyxminiproblems-problemphyxmini0215-durationquantity, def:physics:phyx-mini-0215:phyxminiproblems-problemphyxmini0215-durationreadout}
  Second readout used by the displayed answer choices
\end{definition}
\begin{proof}
  This is the declared model definition of duration In Seconds.
\end{proof}
\begin{definition}[Scene Object]
  \label{def:physics:phyx-mini-0215:phyxminiproblems-problemphyxmini0215-sceneobject}
  \lean{PhyXMiniProblems.ProblemPhyXMini0215.SceneObject}
  Objects whose labels or physical roles are visible in the primary figure
\end{definition}
\begin{proof}
  This is the declared model definition of Scene Object.
\end{proof}
\begin{definition}[Atmospheric Condition]
  \label{def:physics:phyx-mini-0215:phyxminiproblems-problemphyxmini0215-atmosphericcondition}
  \lean{PhyXMiniProblems.ProblemPhyXMini0215.AtmosphericCondition}
  Atmospheric condition explicitly stated in the scenario
\end{definition}
\begin{proof}
  This is the declared model definition of Atmospheric Condition.
\end{proof}
\begin{definition}[Boat Backfire Setup]
  \label{def:physics:phyx-mini-0215:phyxminiproblems-problemphyxmini0215-boatbackfiresetup}
  \lean{PhyXMiniProblems.ProblemPhyXMini0215.BoatBackfireSetup}
  \uses{def:physics:phyx-mini-0215:phyxminiproblems-problemphyxmini0215-lengthquantity, def:physics:phyx-mini-0215:phyxminiproblems-problemphyxmini0215-durationquantity, def:physics:phyx-mini-0215:phyxminiproblems-problemphyxmini0215-speedquantity, def:physics:phyx-mini-0215:phyxminiproblems-problemphyxmini0215-sceneobject, def:physics:phyx-mini-0215:phyxminiproblems-problemphyxmini0215-atmosphericcondition}
  Physical quantities in the boat-to-boat sound-propagation experiment. The directlyAbove relation is an abstract geometric readout from the scene; it records the otherwise calculation-irrelevant placement of boat (b) above tuna school (a) without inventing coordinates not supplied by the image
\end{definition}
\begin{proof}
  This is the declared model definition of Boat Backfire Setup.
\end{proof}
\begin{definition}[Matches Problem And Figure]
  \label{def:physics:phyx-mini-0215:phyxminiproblems-problemphyxmini0215-matchesproblemandfigure}
  \lean{PhyXMiniProblems.ProblemPhyXMini0215.MatchesProblemAndFigure}
  \uses{def:physics:phyx-mini-0215:phyxminiproblems-problemphyxmini0215-lengthinkilometers, def:physics:phyx-mini-0215:phyxminiproblems-problemphyxmini0215-boatbackfiresetup}
  Problem-text and primary-figure readouts. This contains the 1.55 km annotation and the identities of (a), (b), source, and receiver, but no requested elapsed-time value
\end{definition}
\begin{proof}
  This is the declared model definition of Matches Problem And Figure.
\end{proof}
\begin{definition}[Has Physical Parameters]
  \label{def:physics:phyx-mini-0215:phyxminiproblems-problemphyxmini0215-hasphysicalparameters}
  \lean{PhyXMiniProblems.ProblemPhyXMini0215.HasPhysicalParameters}
  \uses{def:physics:phyx-mini-0215:phyxminiproblems-problemphyxmini0215-speedinmeterspersecond, def:physics:phyx-mini-0215:phyxminiproblems-problemphyxmini0215-lengthinmeters, def:physics:phyx-mini-0215:phyxminiproblems-problemphyxmini0215-durationinseconds, def:physics:phyx-mini-0215:phyxminiproblems-problemphyxmini0215-boatbackfiresetup}
  Positivity conditions selecting a physical sound-propagation situation
\end{definition}
\begin{proof}
  This is the declared model definition of Has Physical Parameters.
\end{proof}
\begin{definition}[Uses Standard Air Sound Speed]
  \label{def:physics:phyx-mini-0215:phyxminiproblems-problemphyxmini0215-usesstandardairsoundspeed}
  \lean{PhyXMiniProblems.ProblemPhyXMini0215.UsesStandardAirSoundSpeed}
  \uses{def:physics:phyx-mini-0215:phyxminiproblems-problemphyxmini0215-speedinmeterspersecond, def:physics:phyx-mini-0215:phyxminiproblems-problemphyxmini0215-boatbackfiresetup}
  The standard textbook model input that sound in the ambient air travels at 343 m/s. It is kept separate from both the figure data and the target time
\end{definition}
\begin{proof}
  This is the declared model definition of Uses Standard Air Sound Speed.
\end{proof}
\begin{definition}[Obeys Constant Speed Sound Propagation]
  \label{def:physics:phyx-mini-0215:phyxminiproblems-problemphyxmini0215-obeysconstantspeedsoundpropagation}
  \lean{PhyXMiniProblems.ProblemPhyXMini0215.ObeysConstantSpeedSoundPropagation}
  \uses{def:physics:phyx-mini-0215:phyxminiproblems-problemphyxmini0215-boatbackfiresetup}
  Governing constant-speed propagation law, stated in every unit system through Physlib's dimensionally correct UnitExamples.SpeedEq: v = d / t
\end{definition}
\begin{proof}
  This is the declared model definition of Obeys Constant Speed Sound Propagation.
\end{proof}
\begin{definition}[Answer Choice]
  \label{def:physics:phyx-mini-0215:phyxminiproblems-problemphyxmini0215-answerchoice}
  \lean{PhyXMiniProblems.ProblemPhyXMini0215.AnswerChoice}
  Labels of the four displayed multiple-choice answers
\end{definition}
\begin{proof}
  This is the declared model definition of Answer Choice.
\end{proof}
\begin{definition}[displayed Seconds]
  \label{def:physics:phyx-mini-0215:phyxminiproblems-problemphyxmini0215-answerchoice-displayedseconds}
  \lean{PhyXMiniProblems.ProblemPhyXMini0215.AnswerChoice.displayedSeconds}
  \uses{def:physics:phyx-mini-0215:phyxminiproblems-problemphyxmini0215-answerchoice}
  Number of seconds printed beside each answer label
\end{definition}
\begin{proof}
  This is the declared model definition of displayed Seconds.
\end{proof}
\begin{definition}[Matches Answer Choice]
  \label{def:physics:phyx-mini-0215:phyxminiproblems-problemphyxmini0215-matchesanswerchoice}
  \lean{PhyXMiniProblems.ProblemPhyXMini0215.MatchesAnswerChoice}
  \uses{def:physics:phyx-mini-0215:phyxminiproblems-problemphyxmini0215-answerchoice}
  Agreement with a time displayed to the nearest hundredth of a second
\end{definition}
\begin{proof}
  This is the declared model definition of Matches Answer Choice.
\end{proof}
\begin{lemma}[elapsed Time eq distance div speed]
  \label{lem:physics:phyx-mini-0215:phyxminiproblems-problemphyxmini0215-elapsedtime-eq-distance-div-speed}
  \lean{PhyXMiniProblems.ProblemPhyXMini0215.elapsedTime_eq_distance_div_speed}
  \uses{def:physics:phyx-mini-0215:phyxminiproblems-problemphyxmini0215-speedinmeterspersecond, def:physics:phyx-mini-0215:phyxminiproblems-problemphyxmini0215-lengthinmeters, def:physics:phyx-mini-0215:phyxminiproblems-problemphyxmini0215-durationinseconds, def:physics:phyx-mini-0215:phyxminiproblems-problemphyxmini0215-boatbackfiresetup, def:physics:phyx-mini-0215:phyxminiproblems-problemphyxmini0215-hasphysicalparameters, def:physics:phyx-mini-0215:phyxminiproblems-problemphyxmini0215-obeysconstantspeedsoundpropagation}
  Constant-speed propagation determines elapsed time as distance divided by speed when both are read in SI units. This is a derived relation, not a field of the propagation-law interface
\end{lemma}
\begin{proof}
  The conclusion follows from the governing relations and figure readouts named in the statement of elapsed Time eq distance div speed.
\end{proof}
% --- Archon declaration topology end ---
% --- Archon physics formalization source end ---
